\chapter{Conclusion and Future Directions}
\label{ch:conclusion}

\section{Summary of Contributions}
\label{sec:sum-contrib}

This report has described the design and implementation of
\projname, a single-machine, local-first, conversational AutoML
platform built on the Model Context Protocol. The concrete
contributions are:

\begin{enumerate}[leftmargin=1.4em]
  \item A modular MCP server layout that decomposes a tabular
        pipeline into five independent services and eleven typed
        tools.
  \item An in-process transport layer that removes HTTP overhead
        while preserving the MCP protocol's discovery and
        invocation semantics.
  \item An orchestrator that reduces the LLM to a
        \emph{scheduler}: it never fabricates numbers, only picks
        tools and arguments and formats their results.
  \item A filesystem artefact registry that gives every generated
        file a UUID, a category, and searchable metadata.
  \item A thin Streamlit UI that is a client of the orchestrator
        and can be replaced without touching the ML code.
\end{enumerate}

\section{Findings}
\label{sec:findings}

The end-to-end smoke of the pipeline confirms three important
findings.

\paragraph{The MCP contract is enough.}
A typed tool schema plus a docstring is sufficient information for
current-generation LLMs to plan multi-step workflows. There is no
need for a bespoke protocol on top of MCP for a research-scale
platform.

\paragraph{Deterministic-first works.}
Restricting the LLM to plan-and-format decisions --- and letting
Python do every calculation --- keeps the platform trustworthy. The
user never has to wonder whether a metric was computed or
hallucinated.

\paragraph{Local-first pays off.}
Removing Supabase, authentication, cloud storage, and the Next.js
frontend cut the surface area by more than half and made the
platform tractable for a single-author project without sacrificing
any pipeline capability.

\section{Future Directions}
\label{sec:future}

The following extensions are natural next steps. Each is labelled
\propd{} if a small effort would deliver it, or \future{} for
larger research or engineering programs.

\subsection{AutoML Depth}
\label{sec:future-automl}

\begin{itemize}[leftmargin=1.4em]
  \item \propd{} Add CatBoost \citep{prokhorenkova2018catboost} as a
        fifth bake-off candidate.
  \item \propd{} Add stacking (\code{StackingClassifier} /
        \code{StackingRegressor} \citep{sklearn2011}) as a
        post-bake-off step.
  \item \future{} Integrate deep tabular models (TabNet
        \citep{arik2021tabnet}, FT-Transformer
        \citep{gorishniy2021revisiting}) as optional candidates
        behind a flag.
  \item \future{} Add multi-fidelity HPO (Hyperband
        \citep{li2017hyperband}, BOHB \citep{falkner2018bohb}).
  \item \future{} Feature engineering search
        (Featuretools-style \citep{kanter2015dsm}, LLM-driven CAAFE
        \citep{hollmann2024caafe}).
\end{itemize}

\subsection{Agent Sophistication}
\label{sec:future-agent}

\begin{itemize}[leftmargin=1.4em]
  \item \propd{} Add ReAct-style scratchpad
        \citep{yao2023react} to the system prompt for better
        multi-step reasoning.
  \item \propd{} Add self-consistency
        \citep{wang2023self-consistency} on tool argument
        generation.
  \item \future{} Multi-agent workflows
        \citep{wu2023autogen, hong2024metagpt}: a Planner that
        proposes a DAG of tool calls and a Critic that reviews
        results before final response.
  \item \future{} Retrieval-augmented tool selection over a larger
        catalog \citep{patil2023gorilla, qin2024toolllm}.
\end{itemize}

\subsection{Reproducibility and MLOps}
\label{sec:future-mlops}

\begin{itemize}[leftmargin=1.4em]
  \item \propd{} Record a SHA-256 hash of every uploaded dataset in
        \code{DatasetRecord} and expose it as an artefact metadata
        field.
  \item \propd{} Store leaderboards in a small SQLite database so
        that historical bake-offs can be queried.
  \item \future{} Integrate with MLflow \citep{mlflow} or Weights
        \& Biases \citep{wandb} for experiment tracking and model
        registry semantics.
  \item \future{} Emit OpenTelemetry \citep{opentelemetry} traces
        for every tool call to enable proper observability.
\end{itemize}

\subsection{Deployment and Scaling}
\label{sec:future-deploy}

\begin{itemize}[leftmargin=1.4em]
  \item \propd{} Provide a minimal \file{Dockerfile} that installs
        \file{requirements.txt} and exposes the Streamlit port.
  \item \future{} Split each MCP server into its own process
        (or container) and re-transport them over Streamable HTTP
        \citep{mcpSpec2024} so they can scale independently.
  \item \future{} Add a job queue (Redis / RQ or Celery
        \citep{celery}) so long-running bake-offs do not block the
        UI thread.
  \item \future{} Cloud deployment templates for AWS ECS, Google
        Cloud Run, or Kubernetes.
\end{itemize}

\subsection{Human-in-the-Loop and Guard-rails}
\label{sec:future-hitl}

\begin{itemize}[leftmargin=1.4em]
  \item \propd{} Add a ``dry-run'' mode in which the LLM emits a
        plan of tool calls that the user must approve before
        execution.
  \item \propd{} Add a policy layer that inspects tool arguments
        (e.g.\ reject \code{target\_column} values that do not
        exist).
  \item \future{} Prompt-injection defence: separate tool results
        into a distinct message channel and re-anchor the system
        prompt on every turn.
\end{itemize}

\subsection{Broader Data Modalities}
\label{sec:future-modality}

\begin{itemize}[leftmargin=1.4em]
  \item \future{} Time-series support (statsforecast
        \citep{statsforecast}, sktime \citep{sktime}).
  \item \future{} Text classification with sentence-transformers
        \citep{reimers2019sbert}.
  \item \future{} Image classification via torchvision-style
        pipelines.
  \item \future{} Multimodal LLM-driven captioning of generated
        plots for accessibility.
\end{itemize}

\subsection{Ecosystem Integration}
\label{sec:future-ecosystem}

\begin{itemize}[leftmargin=1.4em]
  \item \propd{} Expose the MCP servers over stdio as well as
        in-process so that Claude Desktop and other MCP hosts can
        use them.
  \item \propd{} Expose MCP resources for the artefact catalog so
        that any MCP client can \code{ReadResource} the leaderboard
        or the champion metadata.
  \item \future{} Publish the MCP servers as a stand-alone package
        that other AutoML front-ends can consume.
\end{itemize}

\subsection{Benchmarking}
\label{sec:future-bench}

\begin{itemize}[leftmargin=1.4em]
  \item \propd{} Automated CI job that runs the protocol of
        Chapter~\ref{ch:experiments} on all six datasets and posts
        the resulting Markdown table as a comment on every PR.
  \item \future{} Compare \projname{} against AutoGluon
        \citep{erickson2020autogluon}, Auto-sklearn
        \citep{feurer2015autosklearn}, and TPOT
        \citep{olson2016tpot} on OpenML CC-18 benchmarks
        \citep{bischl2021openml}.
\end{itemize}

\section{Closing Remarks}
\label{sec:closing}

Conversational, agentic interfaces to data science are neither a
substitute for statistical intuition nor a solution to the
underlying difficulty of applied ML. What they can do, when
carefully designed, is lower the activation energy for someone who
has intuition but has not yet internalised the API surface --- and
give that person a clean, inspectable, downloadable record of
everything the platform did. \projname{} is a small step in this
direction, built to be extended.
